
\subsection{Compromiso de apoyo de la institución}
\label{sec:apoyo}

El proyecto no se desarrolla por encargo de una institución cliente ni cuenta actualmente con una entidad receptora que haya formalizado un compromiso de adopción o implementación del producto. Por esta razón, no se incorpora una carta de compromiso institucional a la propuesta.

No obstante, durante la etapa de identificación de la problemática se ha contado con aportes de profesionales vinculados a la gestión del riesgo y la atención de emergencias. Entre las personas entrevistadas se encuentran:

\begin{itemize}
    \item \textbf{Lina María Marín Rodríguez:} capitán en jefe del sector bomberil, exdirectora de la Dirección Nacional de Bomberos de Colombia y anteriormente directora general del Cuerpo Oficial de Bomberos de Dosquebradas, Risaralda. Su experiencia aporta una perspectiva sobre coordinación institucional, capacidades operativas y atención de incendios forestales \cite{vfpLinaDosquebradas,vfpLinaDNBC2025,vfpLinaDNBC2026}.

    \item \textbf{Ana Milena Mejía Zapata:} directora general de la Defensa Civil Colombiana, entidad dedicada a la gestión del riesgo y la atención de emergencias. Aunque no pertenece a un cuerpo de bomberos, su experiencia en dirección institucional, logística y coordinación de operaciones resulta relevante para comprender las necesidades de información durante una emergencia \cite{vfpAnaDefensaCivil}.

    \item \textbf{Carlos Preciado:} bombero aeronáutico y coordinador regional de servicios de salvamento y extinción de incendios en aeropuertos. Durante la entrevista explicó su responsabilidad en la coordinación de once cuerpos de bomberos aeronáuticos, así como aspectos relacionados con personal, recursos, equipos y capacidad de respuesta \cite{entrevistas2026}.
\end{itemize}

Como contexto normativo de las funciones bomberiles aeronáuticas, el artículo 33 de la Resolución 661 de 2014 contempla a los coordinadores regionales de bomberos dentro de la jerarquía de los Bomberos Aeronáuticos Oficiales de la Aeronáutica Civil, con una homologación de grado correspondiente a capitán-subcomandante \cite{vfpResolucion661SEI}. Esta disposición permite contextualizar institucionalmente el cargo mencionado durante la entrevista, sin que constituya por sí misma una acreditación del grado o de la trayectoria particular del entrevistado.

Las entrevistas han permitido conocer experiencias, necesidades de información y consideraciones operativas relevantes para el planteamiento de la solución \cite{entrevistas2026}.

Estos contactos constituyen un aporte importante para comprender el contexto del problema y podrían facilitar posteriormente la evaluación de la aplicación mediante tareas de consulta e interpretación de escenarios. Su eventual participación en pruebas dependerá de su disponibilidad y aceptación, y no representa un compromiso formal de sus instituciones ni garantiza la adopción del software.

Adicionalmente, se contempla utilizar HPC-ZINE, infraestructura de cómputo de alto desempeño de la Pontificia Universidad Javeriana, como recurso universitario para la preparación de escenarios FDS. La Universidad dispone de recursos de procesamiento, almacenamiento y acompañamiento técnico para proyectos académicos y de investigación \cite{vfpHPCZINE}.

El acceso a este servicio, las condiciones de utilización y los recursos específicos asignados deberán acordarse con sus responsables. Su aprovechamiento como infraestructura de investigación no implica, por sí mismo, que la Universidad o las instituciones de los profesionales entrevistados se constituyan en clientes o receptoras formales de la aplicación.
